% =====================================================================
%  Chapter 18 — The Proof in Natural Language, 2
% =====================================================================

\chapter{The Target: A Positivity Statement}
\label{ch:target}

The first block of the proof is a change of subject, and it is worth doing
carefully because everything afterwards depends on it.

\section{Weil's criterion}For an admissible test function $g$ on the critical strip, write $\hat g$ for
a suitable transform of $g$ --- for instance,
\begin{equation}
\hat g(s) = \int_{\mathbb{R}} g(t)\,e^{-\pi i s t^2}\,dt
\end{equation}
with $g$ smooth, compactly supported and even. Weil's criterion states:

\begin{quote}
\textbf{Proposition B (Weil).} The Riemann hypothesis holds if and only if, for
every admissible $g$,
\[
\sum_{\rho} \hat g(\rho)\,\hat g(1-\rho) \;\ge\; 0,
\]
where the sum is over the non-trivial zeros of $\zeta$, counted with
multiplicity.
\end{quote}

This is a \emph{positivity} condition on a quadratic form in the values
$\hat g(\rho)$. The formalisation of the Weil side is in
\code{Development/WeilBound/}; the criterion itself is stated and proved in
the development's Weil-bound subtree, with the positivity expressed as the
domination of the Weil expression,
\code{Development/WeilBound/ThetaLifts/FrameEstimates/DomWeil.lean}.

\section{Why positivity, and not a bound}The point of the criterion is structural rather than quantitative. A bound on
the sum would give an inequality; a \emph{positivity} gives an exact zero for
the right quantity, and a positivity of a quadratic form can be certified by a
matrix:

\begin{quote}
\textbf{Proposition C.} A real symmetric matrix $M$ is positive semi-definite
if and only if there exist real numbers $a_{ij}$ with
$M = A^{\mathsf T} A$ for some real matrix $A$ --- that is, if and only if
$M$ is a Gram matrix.
\end{quote}

So if one can exhibit the Weil expression as the determinant of a Gram matrix,
RH follows. This is not a technicality; it is the hinge of the whole method.
The remark's item 4 is this sentence:

\begin{quote}
\small
``$W(1S) \ge 0$ via a determinant, interpret as Stepanov with multiplicity one
again; bootstrap!!''
\end{quote}

\section{The Weil expression}The Weil expression $W$ is a real-valued function built from the ordinates of
the zeros. Fix a large real parameter $T$ and write
\begin{equation}
W(T) = \sum_{0 < \gamma \le T} e^{-\pi i \gamma^2 / (2T)} \;+\;
\frac{\gamma(T)}{2} - \int_0^T \psi(x)\,d\!\left(\frac{x}{T}\right),
\label{eq:weilW}
\end{equation}
where $\gamma(T)$ is the number of zeros with imaginary part in $(0, T]$
counted with multiplicity and $\psi(x) = \sum_{n \le x}\Lambda(n)$.

Two features of \eqref{eq:weilW} are used by the argument.

\begin{description}[style=nextline,leftmargin=0pt,labelwidth=0pt,labelsep=0pt,font=\sffamily\bfseries\small,
  itemsep=6pt]
\item[The zeros appear through their ordinates only.] The expression is a sum
of exponentials in $\gamma^2$. It is this that makes the interlacing
structure of the zeros --- the way their arguments alternate --- into a
sign condition on a sum.
\item[The main term is $\int \psi$.] The $\psi$ integral is where the primes
enter, and hence where the arithmetic of the argument has to show up.
\end{description}

The classical statement is that RH is equivalent to $W(T) \ge 0$ for all
sufficiently large $T$ (and, with a little care, for all $T > 0$). The
formalisation's \code{DomWeil} states the domination and \code{FormHeightGen}
and \code{GramMatrixTwistedLieModels} in the HomogeneousNonvanishing subtree
carry the algebra.

\section{The target, restated}\begin{quote}
\textbf{Proposition D.} It suffices to construct, for every admissible $g$, a
matrix $M(g)$ with entries built from automorphic data --- the Gross--Prasad
periods and the interlacing weights of the parameter data --- such that
\[
W(T) \;=\; \det M(g)
\]
and such that $M(g)$ is a Gram matrix.
\end{quote}

This is the reduction that makes the method possible at all. In the
function-field proof of Weil's bound, the analogous Gram matrix is built from
the Frobenius action on the cohomology of the curve: the matrix is the
intersection pairing on a Jacobian, and its positivity is the statement that
Frobenius acts with eigenvalues of absolute value $\sqrt q$. Over the integers
there is no Frobenius and no Jacobian, and the construction must be replaced by
one built from automorphic forms --- which is what ``natively automorphic''
means in the remark.

\section{What is classical and what is not}The formalised development contains the Weil side in substantial form. What it
does \emph{not} contain, in any form anyone outside the author has audited, is
a proof that the matrix $M(g)$ so constructed is a Gram matrix of the right
size and with the right entries. That is the first genuinely new step, and it
is Chapter~\ref{ch:determinant}.

\begin{keytakeaways}
\item Weil's criterion turns RH into a positivity condition on a quadratic
form in the values of a test function at the zeros.
\item A positivity of a quadratic form is certifiable by a Gram matrix, which
is why a determinant is the right object to construct.
\item The Weil expression $W$ is a sum of exponentials in the squares of the
zero ordinates, with a $\psi$-integral main term; the arithmetic of the
argument has to enter through that integral.
\item The whole method reduces to building the Gram matrix out of automorphic
data instead of out of Frobenius on a Jacobian.
\end{keytakeaways}
